\section{The \texttt{Template\_process} McXtrace Component}
Template for a new contributor to create their own physical process.

\subsection*{Identification}
\begin{itemize}
  \item \textbf{Author:} Mads Bertelsen and Erik B Knudsen
  \item \textbf{Origin:} ESS DMSC \& DTU Physics
  \item \textbf{Date:} 20.08.15
\end{itemize}

\subsection*{Description}
This is a template for a new contributor to create their own physical process. The comments in this file are meant to teach the user about creating their own process file, rather than explaining this one. For comments on how this code works, look in the Incoherent\_process.comp.

Part of the Union components, a set of components that work together and thus sperates geometry and physics within McXtrace. The use of this component requires other components to be used.

1) One specifies a number of processes using process components like this one 2) These are gathered into material definitions using Union\_make\_material 3) Geometries are placed using Union\_box / Union\_cylinder, assigned a material 4) A Union\_master component placed after all of the above

Only in step 4 will any simulation happen, and per default all geometries defined before the master, but after the previous will be simulated here.

There is a dedicated manual available for the Union\_components

Algorithm: Described elsewhere

\subsection*{Input parameters}
Parameters in \textbf{boldface} are required; the others are optional.

\begin{longtable}{p{0.22\textwidth}p{0.12\textwidth}p{0.46\textwidth}p{0.14\textwidth}}
\toprule
\textbf{Name} & \textbf{Unit} & \textbf{Description} & \textbf{Default} \\
\midrule
\endhead
sigma & barns & Scattering cross section & 5.08 \\
packing\_factor & 1 & Material packing factor & 1 \\
unit\_cell\_volume & \AA{}$^{3}$ & Unit cell volume & 13.8 \\
interact\_fraction & 1 & How large a part of the scattering events should use this process 0-1 (sum of all processes in material = 1) & -1 \\
\bottomrule
\end{longtable}

\subsection*{Links}
\begin{itemize}
  \item Component source code found in file \texttt{Template\_process.comp}.
\end{itemize}
\IfFileExists{union/Template_process_static.tex}{\input{union/Template_process_static.tex}}{}